\documentclass[9pt,a4paper]{extarticle}

% ======================
% Codificação, idioma e layout
% ======================
\usepackage[T1]{fontenc}
\usepackage[utf8]{inputenc}
\usepackage[brazil]{babel}
\usepackage{geometry}
\geometry{margin=2cm}
\usepackage{parskip}
\usepackage{setspace}
\usepackage{microtype}
\usepackage{enumitem}

% ======================
% Matemática e símbolos
% ======================
\usepackage{amsmath, amssymb, amsthm, bm}

% Algoritmos (CLRS-style)
\usepackage{algpseudocode}

% Cores
\usepackage{xcolor}

% ======================
% Figuras e tabelas
% ======================
\usepackage{graphicx}
\usepackage{booktabs}
\usepackage{array}
\usepackage{tabularx}
\usepackage{float}
\usepackage{caption}
\usepackage{subcaption}

% ======================
% Hiperlinks
% ======================
\usepackage{hyperref}
\hypersetup{
  colorlinks=true,
  linkcolor=blue,
  urlcolor=blue
}

% ======================
% Início do documento
% ======================
\begin{document}

\begin{center}
  \Large \textbf{Prova 1 \\ Introdução à Análise de Algoritmos e Caracterização de Tempos de Execução} \\[4pt]
  \normalsize Disciplina: PCC104 - Projeto de Análise de Algoritmos
\end{center}

\vspace{1cm}

\begin{enumerate}
  \item Por que não é adequado medir a eficiência de um algoritmo contando o número de instruções de máquina executadas em um computador específico? Quais características desejáveis uma boa medida de eficiência deve possuir?

  \item Considere um algoritmo que executa $3^n$ operações e outro que executa $50n^3$ operações. Qual é o aumento do número de operações para cada algoritmo quando o tamanho da entrada é duplicado?

  \item O que significa dizer que um algoritmo é \textit{correto}? Descreva as três propriedades que devem ser verificadas em uma invariante de laço para que ela sirva como ferramenta de prova de corretude.

  \item Considere a seguinte afirmação: \\
  \textit{Se um algoritmo tem complexidade $\Theta(n^2)$, então ele também tem complexidade $O(n^3)$.}\\
  Esta afirmação é verdadeira ou falsa? Explique.

  \item Sejam $f(n)$ e $g(n)$ funções assintoticamente positivas, prove ou refute cada uma das seguintes conjecturas:
  \begin{enumerate}
    \item $f(n) = \Theta(g(n))$ implica $g(n) = \Theta(f(n))$
    \item $\min(f(n), g(n)) = \Omega(f(n) + g(n))$  
    \item $f(n) = O(f(n)^2)$
  \end{enumerate}

  \item Apresente a expressão completa do custo $T(n)$ do algoritmo abaixo:
    \begin{algorithmic}[1]
    \For{$i \gets 1$ \textbf{to} $A.\text{length}$}
        \State $min \gets i$
        \For{$j \gets i + 1$ \textbf{to} $A.\text{length}$}
            \If{$A[j] < A[min]$}
                \State $min \gets j$
            \EndIf
        \EndFor
        \State \textbf{exchange} $A[i]$ \textbf{with} $A[min]$
    \EndFor
  \end{algorithmic}

  \item Apresente a expressão completa do custo $T(n)$ do algoritmo abaixo:
  \begin{algorithmic}[1]
  \For{$i \gets 2$ \textbf{to} $A.\text{length}$}
      \State $key \gets A[i]$
      \State $j \gets 1$
      \While{$j < i$ \textbf{and} $A[j] \leq key$}
          \State $j \gets j + 1$
      \EndWhile
      \For{$k \gets i$ \textbf{downto} $j + 1$}
          \State $A[k] \gets A[k - 1]$
      \EndFor
      \State $A[j] \gets key$
  \EndFor
  \end{algorithmic}

  \item Apresente a expressão completa do custo $T(n)$ do algoritmo abaixo e depois discuta a complexidade de execução deste algoritmo no melhor e no pior caso:
  \begin{algorithmic}[1]
  \Procedure{Find-Max}{$A, n$}
      \State $max \gets A[1]$
      \For{$i \gets 2$ \textbf{to} $n$}
          \If{$A[i] > max$}
              \State $max \gets A[i]$
          \EndIf
      \EndFor
      \State \Return $max$
  \EndProcedure
  \end{algorithmic}
\end{enumerate}

\end{document}
